\SourceLeaf{323}the point sought by 0. I further denote the inclination of the line 12 to a fixed line $L$ simply by 12. This inclination is regarded as increasing in a \textls[120]{specified} sense from 0 to $360^0$ (to fix ideas, suppose $L$ is the meridian and 12 the azimuth of 2 as seen from 1). The signs 21, 13, 31, 23, 32 are to be understood in the same way and referred to the same $L$. It then goes without saying that $21-12=\pm180^0$ (or, more generally, equals an odd multiple of $180^0$). The symbols 01, 02, 03 are to have a similar meaning, with the difference that they are referred, instead of to $L$, to \textls[120]{another entirely arbitrary} fixed line.

One might put it briefly by saying that 12, 13, 21, 23, 31, 32 are the true azimuths of the respective directions, while 01, 02, 03 are the azimuths that have yet to be oriented or that require a correction of orientation whose value is still unknown. It hardly needs saying that azimuth and meridian are used here solely to fix ideas.

The criterion for deciding whether the 6 prescribed quantities (if you count 12 and 21 as only one) are physically compatible is as follows.
\[
\left.
\begin{array}{l}
\sin(12-13-02+03)\\
\sin(23-21-03+01)\\
\sin(31-32-01+02)
\end{array}
\right\}\ \text{must have the same sign.}
\]

The theorem concerning quadrilaterals that I mentioned earlier, which to my knowledge\HistoricalNote{*)}{*)}{Should I be mistaken in this and should you be able to show me that the theorem was already stated by someone else, I would be much obliged to you.} has hitherto been unknown and has been in my possession for nearly half a century, is as follows.
\[
\frac{\sin(12-13-02+03)}{(01).(23)}=\frac{\sin(23-21-03+01)}{(02).(13)}=\frac{\sin(31-32-01+02)}{(03).(12)}\,,
\]

where the factors in the denominators express the distances between the respective points, all taken as positive. The \textls[120]{derivation} of this beautiful theorem by means of complex numbers is of the utmost simplicity. Here I will merely indicate it.



